\uuid{j3oi}
\exo7id{3599}
\titre{exo7 3599}
\auteur{quercia}
\organisation{exo7}
\datecreate{2010-03-10}
\isIndication{false}
\isCorrection{true}
\chapitre{Réduction d'endomorphisme, polynôme annulateur}
\sousChapitre{Autre}
\cpge{mpsi-mp2i.calcul-matriciel-systemes.e, pcsi.determinants.c, ptsi.determinants.a, mpsi-mp2i.matrices-applications-lineaires.a, pcsi.matrices-applications-lineaires.a, ptsi.matrices-applications-lineaires.a, mp-mpi.reduction-endomorphismes.b, pc.reduction.a, psi.reduction-endomorphismes.a, pt.reduction-endomorphismes.a, mp-mpi.reduction-endomorphismes.d, pc.reduction.c, psi.reduction-endomorphismes.c, pt.reduction-endomorphismes.c, pc.complements-espaces-vectoriels.b, psi.complements-espaces-vectoriels.b, pt.complements-algebre-lineaire.b, bcpst2.valeurs-propres-vecteurs-propres.a, bcpst2.valeurs-propres-vecteurs-propres.b}
\module{Algèbre}
\niveau{L2}
\difficulte{}

\contenu{
\question{
Soit $A \in GL_n(\C)$ et $M = \begin{pmatrix} 0 &A \cr I &0 \cr \end{pmatrix} \in \mathcal{M}_{2n}(\C)$.
Montrer que $M$ est diagonalisable si et seulement si $A$ l'est.
(Chercher les sous-espaces propres de $M$ en fonction de ceux de $A$)
}
\reponse{
$E_\lambda(M) = \left\{ \begin{pmatrix}\lambda Y\cr Y\hfil\end{pmatrix} \text{ tq }
         AY = \lambda^2Y \right\}$.
}
}
